\newpage
\begin{center}
\large\textbf{Минобрнауки России}

\large\textbf{Юго-Западный государственный университет}
\vskip 1em
\normalsize{Кафедра программной инженерии}
\vskip 1em
\ifВКР{
        \begin{flushright}
        \begin{tabular}{p{.4\textwidth}}
        \centrow УТВЕРЖДАЮ: \\
        \centrow Заведующий кафедрой \\
        \hrulefill \\
        \setarstrut{\footnotesize}
        \centrow\footnotesize{(подпись, инициалы, фамилия)}\\
        \restorearstrut
        «\underline{\hspace{1cm}}»
        \underline{\hspace{3cm}}
        20\underline{\hspace{1cm}} г.\\
        \end{tabular}
        \end{flushright}
        }\fi
\end{center}
\vspace{1em}
  \begin{center}
  \large
\ifВКР{
ЗАДАНИЕ НА ВЫПУСКНУЮ КВАЛИФИКАЦИОННУЮ РАБОТУ
  ПО ПРОГРАММЕ БАКАЛАВРИАТА}
  \else
ЗАДАНИЕ НА КУРСОВУЮ РАБОТУ (ПРОЕКТ)
\fi
\normalsize
  \end{center}
\vspace{1em}
{\parindent0pt
  Студента \АвторРод, шифр\ \Шифр, группа \Группа
  
1. Тема «\Тема\ \ТемаВтораяСтрока»
\ifВКР{
утверждена приказом ректора ЮЗГУ от \ДатаПриказа\ № \НомерПриказа
}\fi.

2. Срок предоставления работы к защите \СрокПредоставления

3. Исходные данные для создания программной системы:

3.1. Перечень решаемых задач:}

\renewcommand\labelenumi{\theenumi)}

\begin{enumerate}
\item провести анализ предметной области и сравнительный обзор существующих систем биллинга и управления интернет-провайдерами;
\item сформировать функциональные и нефункциональные требования к разрабатываемой системе;
\item спроектировать многоуровневую клиент-серверную архитектуру и структуру реляционной базы данных;
\item разработать программный интерфейс прикладного уровня в стиле REST;
\item реализовать подсистему интеграции с маршрутизаторами MikroTik для синхронизации тарифных ограничений и управления правилами брандмауэра;
\item реализовать подсистему автоматического биллинга: формирование счетов, приостановку услуг абонентам-должникам и автоматическое восстановление обслуживания при погашении задолженности;
\item реализовать подсистему регистрации обращений абонентов сотрудниками технической поддержки;
\item провести модульное и системное тестирование разработанной системы.
\end{enumerate}

{\parindent0pt
  3.2. Входные данные и требуемые результаты для программы:}

\begin{enumerate}
\item Входными данными для программной системы являются: учетные данные
абонентов (фамилия, имя, отчество, документ, удостоверяющий личность,
контакты, адрес установки, координаты, IP-адрес); каталог тарифных планов
(скорости загрузки и отдачи, ежемесячная стоимость, расчетный период);
сведения о платежах и пополнениях лицевых счетов абонентов; параметры
подключения к маршрутизаторам MikroTik по протоколу RouterOS~API; фискальные
реквизиты организации, выпускающей счета.
\item Выходными данными для программной системы являются: сформированные
счета на оплату с печатной формой в формате PDF; синхронизированные простые
очереди и правила брандмауэра на маршрутизаторе MikroTik; актуальные статусы
обслуживания абонентов; журнал аудита изменений; аналитические показатели
деятельности оператора (численность активных и приостановленных абонентов,
объем выставленных и оплаченных средств, перечень крупнейших дебиторов).
\end{enumerate}

{\parindent0pt

  4. Содержание работы (по разделам):
  
  4.1. Введение.

  4.2. Анализ предметной области: бизнес-процессы интернет-провайдера,
маршрутизаторы MikroTik и протокол RouterOS~API, сравнительный анализ
существующих решений.

4.3. Техническое задание: основание для разработки, цель и назначение
разработки, требования к функциональности и интерфейсу, моделирование
вариантов использования, требования к надежности, информационной безопасности
и совместимости.

4.4. Технический проект: общая характеристика проектного решения, выбор
средств и технологий реализации, архитектурное построение системы, проект
базы данных, подсистема интеграции с маршрутизаторами MikroTik, подсистема
автоматизации.

4.5. Рабочий проект: структура серверной и клиентской частей, тестирование
разработанной системы.

4.6. Заключение.

4.7. Список использованных источников.

5. Перечень графического материала:

Лист 1. Сведения о ВКРБ.

Лист 2. Цель и задачи разработки.

Лист 3. Диаграмма прецедентов.

Лист 4. Архитектура системы.

Лист 5. ER-модель базы данных.

Лист 6. Реляционная модель базы данных.

Лист 7. Заключение.

\vskip 2em
\begin{tabular}{p{6.8cm}C{3.8cm}C{4.8cm}}
Руководитель \ifВКР{ВКР}\else работы (проекта) \fi & \lhrulefill{\fill} & \fillcenter\Руководитель\\
\setarstrut{\footnotesize}
& \footnotesize{(подпись, дата)} & \footnotesize{(инициалы, фамилия)}\\
\restorearstrut
Задание принял к исполнению & \lhrulefill{\fill} & \fillcenter\Автор\\
\setarstrut{\footnotesize}
& \footnotesize{(подпись, дата)} & \footnotesize{(инициалы, фамилия)}\\
\restorearstrut
\end{tabular}
}

\renewcommand\labelenumi{\theenumi.}
